% =============================================================
% Conclusion générale et perspectives — version compacte (~3 pages)
% =============================================================

\documentclass[14pt,a4paper,openany]{extreport}

\usepackage[T1]{fontenc}
\usepackage[utf8]{inputenc}
\usepackage[french]{babel}
\addto\captionsfrench{%
  \renewcommand{\tablename}{Tableau}%
  \renewcommand{\figurename}{Figure}%
}

\usepackage{lmodern}
\usepackage{setspace}
\onehalfspacing
\usepackage[a4paper,inner=3cm,outer=2.3cm,top=2.5cm,bottom=2.7cm,headheight=15pt,bindingoffset=0.3cm]{geometry}

\usepackage{amsmath,amssymb}
\usepackage{xcolor}
\definecolor{uridarkblue}{RGB}{30,40,90}
\definecolor{urigray}{RGB}{70,70,70}
\definecolor{uriaccent}{RGB}{180,60,40}

\usepackage{array}
\usepackage{booktabs}
\usepackage{tabularx}
\usepackage[most]{tcolorbox}

\usepackage{titlesec}
\titleformat{\section}
  {\normalfont\Large\bfseries\color{uridarkblue}}
  {\thesection.}{0.8em}{}

\usepackage{fancyhdr}
\pagestyle{fancy}
\fancyhf{}
\fancyhead[L]{\textsc{\nouppercase{\leftmark}}}
\fancyhead[R]{\thepage}
\fancyfoot[L]{\footnotesize Mémoire de Master en Informatique, Université de Dschang}
\fancyfoot[R]{\footnotesize URIFIA}
\renewcommand{\headrulewidth}{0.4pt}
\renewcommand{\footrulewidth}{0.2pt}

\usepackage[font=small,labelfont={sc,bf},textfont=it]{caption}
\usepackage[numbers,square,sort&compress]{natbib}
\usepackage[hidelinks]{hyperref}

\newcommand{\Oracle}{\textsc{Oracle}}
\newcommand{\CorrDiff}{\textsc{CorrDiff}}
\newcommand{\HR}{\ensuremath{\text{HR}}}
\newcommand{\LR}{\ensuremath{\text{LR}}}
\newcommand{\muHR}{\ensuremath{\mu_{\HR}}}
\newcommand{\Adag}{\ensuremath{A_{\mathrm{dag}}}}

\setcounter{page}{100}

\begin{document}

\chapter*{Conclusion générale et perspectives}
\addcontentsline{toc}{chapter}{Conclusion générale et perspectives}
\markboth{Conclusion générale et perspectives}{Conclusion générale et perspectives}

% =============================================================
\section*{Problématique étudiée et choix méthodologiques}
\label{concl:probl}
% =============================================================

Ce mémoire est parti d'un constat : aucun downscaler génératif probabiliste actuel ne combine la rapidité des approches profondes avec une robustesse de principe au changement de distribution. Le paradigme à deux étages de \CorrDiff{}~\cite{mardani_corrdiff_2023} stabilise l'entraînement mais son premier étage de régression générique reste vulnérable à tout déplacement de régime. Notre choix méthodologique a été de \emph{causaliser} ce premier étage en plaçant la causalité dans la topologie du graphe de calcul plutôt que dans une régularisation de la perte --- une régularisation peut être contournée par un réseau expressif, une dépendance architecturale est par construction indéfaisable.

Quatre éléments structurent la contribution : (\textbf{C1}) une architecture two-stage où la causalité est inscrite dans la topologie du chemin de calcul ; (\textbf{C2}) un réseau causal récurrent (RCN) intégrant SCM dynamique, GRU~\cite{cho_gru_2014} et pénalité DAGMA~\cite{bello_dagma_2022} dans une seule cellule ; (\textbf{C3}) une interface inter-étages par concaténation de canaux garantissant qu'aucune information ne contourne le chemin causal (propriété O3 vérifiée empiriquement, ratio $0{,}87$) ; (\textbf{C4}) un protocole expérimental à six familles de métriques avec deux diagnostics causaux originaux --- ratio O3 et test d'intervention $Q_{\mathrm{int}}$.

% =============================================================
\section*{Analyse critique des résultats obtenus}
\label{concl:analyse}
% =============================================================

Sur la Nouvelle-Zélande (ACCESS-CM2 en distribution ; EC-Earth3 et NorESM2-MM en inter-GCM historique), \Oracle{} réduit le CRPS de $40\,\%$ en distribution et de $30\,\%$ en inter-GCM par rapport à une baseline non-causale strictement comparable. Le ratio \emph{spread-skill} passe de $0{,}20$ à $0{,}69$, soit un facteur $3{,}5$ d'amélioration de la calibration d'ensemble. La distance spectrale RAPSD chute d'un facteur $450$. Sur deux perturbations contrôlées, \Oracle{} répond avec un signe physiquement cohérent ($Q_{\mathrm{int}}=2/2$), là où la baseline non-causale échoue sur l'intervention thermique. L'hypothèse H2 (séparation moyenne déterministe + résidu stochastique $\to$ stabilité et calibration) est validée sans ambiguïté ; H1 (causalité $\to$ généralisation inter-GCM) est validée modérément ; H3 (O3 architectural $\to$ pas de DAG décoratif) est validée par construction.

Plusieurs limites doivent cependant être assumées explicitement :
\begin{itemize}[leftmargin=1.2em, itemsep=2pt]
    \item La matrice $\Adag$ apprise reste à magnitudes uniformes ($Q_{\mathrm{phys}} = 0{,}40$, équivalente à un graphe tiré dans la classe d'équivalence de Markov), conformément à la limite théorique d'identifiabilité sur données observationnelles~\cite{pearl_causality_2009} : le DAG appris est une régularisation structurelle inductive, pas une cartographie causale identifiée.
    \item Les magnitudes interventionnelles restent environ trois ordres de grandeur sous Clausius-Clapeyron théorique --- le signe est correct, la magnitude ne l'est pas.
    \item Le test inter-GCM porte sur le régime historique, distinct d'un vrai test hors-distribution climatique (SSP futur à effectuer).
    \item Aucune donnée africaine n'est testée et la métrique F1@p99 régresse en raison d'une régularisation interne lissante.
    \item Les chiffres reposent sur un seul seed d'entraînement (contrainte CPU) et une taille d'ensemble modeste ($K=12$ inter-GCM, $K=64$ ID, contre 192~membres pour \CorrDiff{}).
\end{itemize}

% =============================================================
\section*{Quelques perspectives}
\label{concl:persp}
% =============================================================

Cinq axes structurent la continuation, hiérarchisés par horizon temporel :

\textbf{Court terme (1--2 mois).} Recalibration d'ensemble post-hoc par EMOS~\cite{gneiting_emos_2005} ou BMA, écartée ici pour préserver la comparabilité avec la baseline. Validation sur SSP5-8.5 : exécutable sans modification architecturale sur des données CMIP6 disponibles, ce qui transformerait un test inter-GCM historique en véritable test hors-distribution climatique.

\textbf{Moyen terme (3--6 mois).} Ancrage prédictif joint de type CASTLE~\cite{kyono_castle_2020}, forçant chaque arête de $\Adag$ à porter un signal effectif par \emph{fine-tuning} de 25--30 époques sans changement d'architecture. Masque physique sur $\Adag$ via un terme $\lambda_{\mathrm{prior}} \Vert \Adag - \alpha G_{\mathrm{phys}} \Vert^2$ pénalisant les écarts au couplage quasi-géostrophique attendu, dans la lignée du \emph{physics-informed ML}~\cite{karniadakis_pinn_2021}.

\textbf{Long terme (1--3 ans).} Transfert vers l'Afrique subsaharienne, motivation finale de ce travail : \emph{sanity-check} zéro-shot sur les Hautes Terres éthiopiennes ou le Sahel ouest avec CHIRPS comme cible HR, puis ré-entraînement sur prédicteurs adaptés (mousson, ondes d'Est, ITCZ) et DAG révisé --- le prior NZ encode un couplage descendant (jet polaire), alors qu'en Afrique équatoriale la convection profonde impose un schéma ascendant (IVT~$\to$~CAPE~$\to$~SP$_{\mathrm{HR}}$), avec ERA5~\cite{hersbach_era5_2020} comme proxy d'initialisation faute de réseau radar dense.

Ce mémoire défend que la causalité utile au downscaling probabiliste n'est pas une propriété décorative émergeant d'une régularisation, mais une contrainte architecturale à placer explicitement dans la topologie du graphe de calcul --- en Nouvelle-Zélande aujourd'hui, dans les communautés d'Afrique centrale et de l'Ouest qui ont motivé ce travail demain.

\end{document}
